% Opening: the title and Doge, to the left of the map; then the camera pulls back to the map.

\gmgo{title}{MemeAtlas}
{\setbeamertemplate{background}{}
\begin{frame}[plain]
  \begin{tikzpicture}[remember picture, overlay]
    \coordinate (c) at ($(current page.center)+(4.3cm,-.1cm)$);
    \foreach \a/\r/\col in {15/3.05/kWikidata, 48/2.55/kTemplate, 88/2.95/kFrame, 128/2.6/kEvent,
                            205/3.0/kConcept, 248/2.55/kWikidata, 292/2.95/kEvent, 328/2.5/kFrame} {
      \draw[gmBase, line width=.8pt] (c) -- ($(c)+(\a:\r)$);
      \fill[\col] ($(c)+(\a:\r)$) circle (.2);}
    \node[circle, minimum size=4.3cm, draw=kFrame, line width=3pt,
          path picture={\node at (path picture bounding box.center)
                        {\includegraphics[height=4.5cm]{figures/img/doge}};}] at (c) {};
    \node[anchor=south west, font=\fontsize{44}{48}\selectfont\bfseries, text=gmInk]
      at ($(current page.west)+(11mm,4mm)$) {MemeAtlas};
    \node[anchor=north west, text width=7.4cm, font=\Large, text=gmInk2, align=left]
      at ($(current page.west)+(11.5mm,1mm)$) {\hyphenpenalty=10000 Turning Know Your Meme into a living knowledge graph};
    \node[anchor=south west, font=\large, text=gmInk2, align=left]
      at ($(current page.south west)+(11.5mm,9mm)$) {\insertauthor\\[1mm]{\normalsize\insertinstitute\quad\textperiodcentered\quad\insertdate}};
  \end{tikzpicture}
  \note{
  \begin{itemize}
    \item Welcome. One hour, one meme, eleven stops.
    \item The picture already says the whole talk: a meme in the middle, and everything we
      know about it connected around it, each colour a different kind of knowledge.
    \item Hold on to that dog. We will follow him through every step.
  \end{itemize}}
\end{frame}}

\gmgo{meet}{Meet Doge}
\begin{frame}{Meet Doge}
  \begin{columns}[c]
    \column{.42\textwidth}\centering
      \begin{tikzpicture}
        \node[circle, minimum size=4.3cm, draw=kFrame, line width=3pt,
              path picture={\node at (path picture bounding box.center)
                            {\includegraphics[height=4.5cm]{figures/img/doge}};}] {};
      \end{tikzpicture}\par\vskip2mm
      {\normalsize\color{gmMuted}Kabosu, 2010 --- Know Your Meme's ``Doge''}
    \column{.54\textwidth}
      \LARGE
      \begin{itemize}\setlength{\itemsep}{4mm}
        \item Who is this dog?
        \item Where did the joke start?
        \item Which memes grew out of it?
      \end{itemize}
      \vskip7mm
      \Large\color{gmInk2}A computer can answer only if the facts are \textbf{\color{gmInk}connected}.
  \end{columns}
  \note{
  \begin{itemize}
    \item Ask the room: who knows this dog? (Everyone does.)
    \item Now ask a computer. It can find pages that mention Doge, but it cannot \emph{answer}:
      it does not know that the dog is a Shiba Inu named Kabosu, that the joke started on
      Tumblr in 2010, that Dogecoin and Cheems came out of it.
    \item People know it because the facts are connected in our heads. My project connects them
      for the computer: a knowledge graph of memes.
  \end{itemize}}
\end{frame}

\gmoverview{One meme,\\ eleven stops}
\note{
\begin{itemize}
  \item This is the talk, and it is also the system: each circle is one step of the pipeline,
    each arrow is data flowing from one step to the next. Everything we will see lives inside
    one of these circles.
  \item Grey, top left: where the project started, IMKG. The dashed box: the machine that runs
    every step, every month. Pink, at the top: the vocabulary the graph is written in.
  \item Doge (the small photo) walks the map with us. In each circle: first what happens to him
    there, then the engineering decisions behind it.
\end{itemize}}
